% from main.tex
%%%%%%%%%%  scan idx 695  |  folio 682  %%%%%%%%%%
mais tels qu'on ait
\[
 c_1(L_1)c_1(L_2)c_1(L_3)\ne0
 \quad\text{dans }H^6(X,\mu_2^{\otimes3}).
\]

\textbf{5.3.} Lorsque $X$ est lisse et quasi-projectif sur un corps $k$, on peut de même se poser la question de l'existence d'un homomorphisme d'anneaux
\[
 g:\operatorname{Gr}^\bullet(X)\longrightarrow A(X)
\]
à valeurs dans l'anneau de Chow $A(X)$, transformant classes de Chern en classes de Chern, et on peut répéter à ce propos les reflexions qui précèdent, qui montrent en particulier comment construire directement $(-1)^{i-1}(i-1)!g^i$ à défaut de $g^i$ lui-même.

\textbf{5.4.} Signalons, pour terminer ce numéro, qu'essentiellement la même question se pose si, au lieu d'un schéma $X$, on considère un espace topologique compact (disons) $X$, ce qui permet encore de considérer le gradué $\operatorname{Gr}^\bullet(X)$ de $K^\bullet(X)$ muni de sa $\lambda$-filtration, et de se proposer de trouver un homomorphisme d'anneaux
\[
 \operatorname{Gr}^\bullet(X)\longrightarrow H^{2\bullet}(X,\Z)
\]
transformant classes de Chern en classes de Chern.

\section*{6. Théorème de Riemann--Roch cohomologique, et homomorphisme de Gysin}

\textbf{6.1.} Pour simplifier, nous nous bornerons ici aux schémas $X$ admettant un Module inversible ample, ce qui implique que l'on n'a pas à distinguer entre le $K'(X)$ naif et l'autre. Commençons par la formule de Riemann--Roch ``avec dénominateurs'', pour un morphisme propre $f:X\to Y$ d'intersection complète. Comme (pour un nombre premier $\ell$ premier aux caractéristiques résiduelles de $Y$) on peut définir les classes de Chern $\ell$-adiques d'un Module localement libre $E$,
\[
 c_i(E)\in H^{2i}(X,\Zell(i))
\]
(SGA 5 VII), correspondant à un homomorphisme
\begin{equation}\tag{6.1}
 K'(X)\longrightarrow \Ch\bigl(H^{2\bullet}(X,\Zell(\bullet))\bigr),
\end{equation}
la question se pose de donner une variante cohomologique de la formule de Riemann--Roch. Le seul ingrédient qui manque pour donner un sens à la formule est un \emph{homomorphisme de Gysin}
